\section{Technical Research Lines}
\label{sec:methods}
A useful taxonomy has at least two independent layers: the graphic-design problem being solved and the computational mechanism used.
The same representation may be produced by diffusion, autoregressive models, optimization, or multi-stage agents, while the same mechanism may produce several kinds of artifact.

\subsection{Constraint-Aware Layout Generation}
Layout generation predicts discrete elements and their geometric attributes.
This includes unconstrained synthesis, category-conditioned generation, partial completion, and constraint-aware refinement.
CGL-GAN conditions advertising layouts on visual context \citep{cglgan2022}; LayoutDETR applies detection-transformer-style conditioning to foreground composition \citep{layoutdetr2024}.
LayoutDM and LayoutDiffusion illustrate diffusion over discrete layout attributes \citep{layoutdm2023,layoutdiffusion2023}, while LayoutFlow addresses continuous geometry \citep{layoutflow2024}.
Early and intermediate directions include wireframe discrimination in LayoutGAN \citep{layoutgan2019}, VAE-based document reconstruction in CanvasVAE \citep{canvasvae2021}, self-attentive scene completion in LayoutTransformer \citep{layouttransformer2021}, latent-space constraint optimization \citep{constraintlayout2021}, diffusion with explicit aesthetic constraints in LACE \citep{lace2024}, and in-context layout generation with LayoutPrompter \citep{layoutprompter2023}.
LayoutGPT uses language models for visual planning across domains \citep{layoutgpt2023}.
PosterLayout supplies a benchmark and method for content-aware visual-textual arrangement \citep{posterlayout2023}.
Importantly, the name ``LayoutDiffusion'' is also used for a distinct layout-to-image synthesis method \citep{layoutdiffusion_l2i2023}; its image-generation results should not be mixed with the discrete graphic-layout model \citep{layoutdiffusion2023}.
In comparing these systems, the allowed conditions and datasets must be matched: unconditional realism, completion, and relational constraint satisfaction are not equivalent tasks.


\paragraph{Conditional layout families.}
LayoutVAE samples arrangements from object-label sets rather than solving fully specified graphic-design tasks \citep{layoutvae2019}, whereas BLT iteratively completes masked attributes with a bidirectional transformer \citep{blt2022}.
LayoutFormer++ serializes heterogeneous geometric constraints and restricts decoding \citep{layoutformerpp2023}, while DLT models discrete categories and continuous coordinates jointly \citep{dlt2023}.
Dolfin studies diffusion layout transformers without a separate autoencoder \citep{dolfin2023}, whereas FlexDM predicts masked fields including text, styling, and images in multimodal design documents \citep{flexdm2023}.
CoLay supports combinations of text, guidelines, element categories, and partial layout conditions \citep{colay2024}, while Layout-Corrector repairs defective diffusion samples in the generation process \citep{layoutcorrector2024}.
LayoutRectifier uses a separate geometric optimization pass to reduce alignment, overlap, and containment errors \citep{layoutrectifier2025}, whereas StructLayoutFormer generates explicit hierarchical layout structures \citep{structlayoutformer2025}.
These are distinct condition contracts, not interchangeable measurements of controllable design.

\paragraph{Code, retrieval, and spatial reasoning.}
LayoutNUWA formulates layout completion as HTML code generation \citep{layoutnuwa2024}, while PosterLlama applies code-formatted layouts and augmentation to content-aware poster design \citep{posterllama2024}.
RALF retrieves similar designs to guide image-conditioned generation \citep{ralf2024}, and RADM jointly models visual-textual and geometric relations \citep{radm2023}.
ReLayout reasons about content regions, margins, and element relations for conditional placement \citep{relayoutreason2025}, while LaySPA adds reward-guided spatial reasoning to language-model layout planning \citep{layspa2025}.
Uni-Layout combines generation and human-feedback-aligned assessment \citep{unilayout2025}, whereas iPoster emphasizes interactive specification of partial constraints \citep{iposter2026}.
Scan-and-Print instead investigates efficient patch selection and training augmentation for content-aware poster layouts \citep{scanprint2025}.
These approaches require matched input constraints and output granularity before their geometric metrics can support a comparison.

\subsection{Multi-Stage Graphic Design and Agents}
Integrated poster generation must coordinate language, asset production, layout decisions, and visual verification.
COLE and OpenCOLE use hierarchical design decomposition \citep{cole2023,opencole2024}.
CreatiDesign develops multi-conditional visual generation across primary content and decorative or textual components \citep{creatidesign2026}.
Vinci combines data-driven poster creation with interactive editing feedback and user evaluations \citep{vinci2021}.
AutoPoster extends content-aware ad generation through layout, taglines, and predicted style attributes \citep{autoposter2023}; Graphist infers ordering and spatial arrangements over unordered input assets \citep{graphist2025}; and BannerAgency produces component-based banners from brand requirements \citep{banneragency2025}.
PosterCraft addresses raster poster generation with text-focused and aesthetic objectives \citep{postercraft2026}.
Code-oriented systems increasingly produce the design document itself, as in DesignAsCode and Editable Visual Design \citep{designascode2026,editablevisual2026}; PSDesigner targets layered editor workflows \citep{psdesigner2026}.
PosterAgent studies staged draft generation and refinement through agentic reinforcement learning \citep{posteragent2026}.
Visual-feedback refinement and OCR-aware reinforcement learning are investigated in VFLM \citep{guo2026visualfeedback}.
We separate \emph{reflection during generation} from an edit request received after a completed artifact is delivered.


\paragraph{Commercial, personalized, and artistic workflows.}
T-Stars-Poster combines prompt preparation, layout, background generation, and rendering for product advertising \citep{tstarsposter2025}, while reliable ad generation uses feedback-driven inspection to identify unsuitable outputs \citep{reliablead2024}.
POSTA separates background, layout, and rendering for customized artistic posters \citep{posta2025}, whereas PosterVerse develops a commercial poster workflow with HTML-oriented text rendering \citep{posterverse2026}.
SimplePoster studies product preservation and constrained text rendering through image inpainting \citep{simpleposter2026}, while InnoAds-Composer jointly conditions product content, text, and style \citep{innoads2026}.
Design Your Ad models personalized image-and-text advertising based on user histories \citep{designyourad2026}, and AutoPP introduces poster optimization informed by online feedback \citep{autopp2025}.
These signals target different qualities: engagement, text correctness, product identity, and human aesthetic preference cannot be conflated.
Brief2Design supports requirements clarification, element exploration, and recombination from professional briefs rather than treating the completed image as the only useful output \citep{brief2design2026}.

\subsection{Typography and Vector Synthesis}
Text rendering poses constraints not captured by photographic image fidelity.
TextDiffuser and TextDiffuser-2 make text layout a first-class part of image generation \citep{textdiffuser2023,textdiffuser22024}; GlyphDraw2 targets stylized, multilingual glyph-aware poster generation \citep{glyphdraw2025}; GlyphControl and AnyText provide complementary glyph-conditioned and multilingual visual text generation/editing approaches \citep{glyphcontrol2023,anytext2024}.
TextPainter conditions poster typography on both local visual context and language semantics \citep{textpainter2023}.
Word-As-Image optimizes letter shapes to communicate meaning while retaining legibility and font resemblance \citep{wordasimage2023}.
De-Rendering Stylized Texts instead infers editable rendering parameters from a flattened text image \citep{derendering2021}.
A replacement of semantic text, a modification of glyph outlines, and a reconstruction of font/effect parameters are three different editing operations.

Vector-centric work generates or embeds shape/path sequences.
DeepSVG and SVGformer learn vector representations \citep{deepsvg2020,svgformer2023}, whereas VectorFusion and IconShop synthesize vectors conditioned on textual prompts \citep{vectorfusion2023,iconshop2023}.
Neural-path generation explicitly models shape structure \citep{neuralpath2024}, while LayerTracer synthesizes layered SVG designs through staged diffusion and vectorization \citep{layertracer2025}.
Code-generation approaches such as StarVector and Chat2SVG explore additional SVG primitives and text-conditioned editing workflows \citep{starvector2025,chat2svg2025}.
Their visual or semantic generation scores should not be interpreted as evidence for the broader class of targeted design edits.

\subsection{Reconstruction and Explicit Editing}
Another line starts from an existing artifact.
LayerD decomposes rendered designs into layers \citep{layerd2025}; Accordion generates layered designs by planning and decomposing a complete visual reference \citep{accordion2025}; PosterCopilot reasons about composition and repeated layer-sensitive changes \citep{postercopilot2026}.
MiLDEdit explicitly addresses instruction-driven modifications of multi-layer design documents, presenting a dedicated benchmark and evaluation framework \citep{mildedit2026}.
PosterOmni spans local raster-image editing and global poster creation \citep{posteromni2026}.
These tasks are meaningfully editing-oriented, but a pixel-space editing benchmark does not establish that the result is an editable PSD or DOM structure.
Each study must be coded against its actual inputs, outputs, and scoring procedure.

\subsection{Slides, Infographics, and Scientific Figures}
Presentation slides form a distinct task family rather than a synonym for a poster.
AutoPresent generates structured slide visuals from natural-language specifications and compares image-level generation with programmatic outputs \citep{autopresent2025}.
PPTAgent converts source documents into presentations using reference-based planning and editing actions, and evaluates content, design and cross-slide coherence \citep{pptagent2025}.
This distinction matters: the visual fidelity of one slide and the logical coherence of an entire deck are separate observables.

Infographics also demand more than pleasing composition.
BizGen studies article-level visual text rendering in densely structured infographics, where correct long-form text and layout binding are central \citep{bizgen2025}.
InfoEdit instead starts from an existing information graphic and evaluates editing instructions that can require neighboring components to reflow \citep{yang2026infoedit}.
The former is a generation problem; the latter tests post-generation editing.
Their metrics and output contracts should not be pooled.

Scientific figures and academic posters convey factual research content, introducing constraints on semantic correctness and reading order.
AutoFigure creates scientific illustrations from extended source material \citep{autofigure2026}, while PosterGen creates research posters through content selection, layout, typography and rendering agents \citep{postergen2026}.
The shared R--E--V axes remain useful across domains, but comparisons must respect whether the delivered artifact is a vector, a presentation document, or merely a rendered image.
